\documentclass[11pt]{article}
\usepackage[margin=1in]{geometry}
\usepackage{hyperref}
\usepackage{xcolor}

\hypersetup{
    colorlinks=true,
    linkcolor=blue,
    urlcolor=blue,
}

\title{Assignment 2: Ethereum --- Phase 2 Submission}
\author{Blockchains and Cryptocurrencies (601.641/441) --- Spring 2026}
\date{}

\begin{document}
\maketitle

\section*{Student Information}

\begin{tabular}{l l}
\textbf{Name (string):} & \texttt{TODO: Firstname Lastname} \\[6pt]
\textbf{JHED ID:}        & \texttt{TODO: abc123} \\
\end{tabular}

\section*{On-Chain Information}

\begin{tabular}{l p{10cm}}
\textbf{Wallet Address:}       & \texttt{TODO: 0x...} \\[6pt]
\textbf{StudentToken Address:} & \texttt{TODO: 0x...} \\[6pt]
\textbf{SimplePair Address:}   & \texttt{TODO: 0x...} \\[6pt]
\textbf{Swap Tx Hash:}         & \texttt{TODO: 0x...} \\
\end{tabular}

\vspace{1em}

\noindent\textbf{Note:} The TA will verify your submission by:
\begin{enumerate}
    \item Checking that your SimplePair contract is deployed and initialized with your StudentToken and the GOV token
    \item Verifying the swap transaction executed successfully through your pair
    \item Running the autograder against your \texttt{SimplePair.sol} source code
\end{enumerate}

\vspace{1em}

\section*{Short Answer}

\noindent\textbf{Q1:} The first liquidity provider receives $\sqrt{amountA \times amountB}$ LP tokens.

\begin{itemize}
    \item[(a)] Alice deposits 100 tokenA and 400 tokenB into an empty pool. How many LP tokens does she receive?
    \item[(b)] Now suppose we used $amountA + amountB$ instead. Alice could deposit 1 tokenA and 499 tokenB (total = 500 LP tokens) instead of 100 and 400 (total = 500 LP tokens). Explain concretely what problem this creates for the pool and for future swappers.
    \item[(c)] How does the geometric mean ($\sqrt{amountA \times amountB}$) fix this problem? What happens to Alice's LP token count if she tries the 1/499 split with the sqrt formula?
\end{itemize}

\vspace{5em}

\noindent\textbf{Q2:} A pool has 10{,}000 tokenA and 10{,}000 tokenB. You swap 1{,}000 tokenA in. How many tokenB do you receive? Show your work using the constant product formula with the 0.3\% fee.

\vspace{4em}

\noindent\textbf{Q3:} Consider three AMM invariants:
\begin{itemize}
    \item \textbf{Constant sum:} $x + y = k$ \quad (always swaps 1:1)
    \item \textbf{Constant product:} $x \cdot y = k$ \quad (Uniswap V2)
    \item \textbf{StableSwap (Curve):} a hybrid that behaves like constant sum near the midpoint and like constant product at the extremes
\end{itemize}

\begin{itemize}
    \item[(a)] A constant sum pool holds 10{,}000 USDC and 10{,}000 DAI. An attacker swaps in 10{,}000 DAI. What does the attacker receive? What is left in the pool? Explain why this makes the constant sum model broken as a general-purpose AMM.
    \item[(b)] Now consider the same trade on a constant product pool ($k = 10{,}000 \times 10{,}000 = 10^8$). The attacker swaps in 10{,}000 DAI. Calculate how much USDC they receive (ignore fees). Why does the constant product model survive this attack?
    \item[(c)] A USDC/USDT pool on Curve uses StableSwap. A user swaps 1{,}000 USDC for USDT. On Uniswap V2 (constant product) with the same reserves, the same swap yields fewer USDT due to slippage. Why does StableSwap give a better rate for stablecoin pairs, and why would it be \textit{dangerous} to use StableSwap for a volatile pair like ETH/USDC?
\end{itemize}

\vspace{5em}

\noindent\textbf{Q4:} In the \texttt{SimpleFactory} contract, calling \texttt{createPair()} deploys a new \texttt{SimplePair} contract from within the factory.

\begin{itemize}
    \item[(a)] What EVM opcode is used when a smart contract deploys another smart contract using the \texttt{new} keyword in Solidity? How does this differ from the \texttt{CREATE2} opcode?
    \item[(b)] Why would a protocol like Uniswap prefer \texttt{CREATE2} over \texttt{CREATE} for deploying pair contracts? \textit{Hint: think about what determines the deployed address in each case.}
\end{itemize}

\vspace{5em}

\noindent\textbf{Q5:} \textit{Impermanent Loss} (IL).

You deposit 1 ETH and 2{,}000 USDC into a constant product pool (ETH = \$2{,}000). The pool now has reserves $x = 1$ ETH and $y = 2{,}000$ USDC, so $k = 2{,}000$.

The price of ETH rises to \$4{,}000. Arbitrageurs trade with the pool until the pool price matches the market price.

\begin{itemize}
    \item[(a)] After arbitrage, the pool price must satisfy $\frac{y}{x} = 4{,}000$ and $x \cdot y = 2{,}000$. Solve for the new reserves $x$ and $y$. \textit{(Show your algebra.)}
    \item[(b)] Calculate the total USD value of your LP position using the new reserves from (a). Then calculate what your tokens would be worth if you had simply held 1 ETH + 2{,}000 USDC in your wallet. What is the impermanent loss in dollars and as a percentage?
    \item[(c)] You earned 50 USDC in trading fees while providing liquidity. Does this fully offset the impermanent loss from (b)?
    \item[(d)] You now withdraw your liquidity. Is the loss still ``impermanent''? Explain in 1--2 sentences when impermanent loss becomes a realized (permanent) loss.
\end{itemize}

\vspace{5em}

\noindent\textbf{Q6:} Explain in 2--3 sentences why an AMM user cannot extract more of a token from the pool than was deposited into it.

\vspace{4em}

\end{document}
